\begin{exercise}{12}
Show that for every simplicial set $X$, there is a canonical bijection
$\pi_0^\Delta(X)\cong\pi_0(|X|)$.

\begin{proof}
Sending a vertex to its point in $|X|$ induces a map
$\phi\colon\pi_0^\Delta(X)\to\pi_0(|X|)$: the endpoints of each
edge are joined by its realized path. Every point of $|X|$ is
represented by a point in a topological simplex, and a straight-line
path in that simplex joins it to a vertex. Hence $\phi$ is surjective.

To prove injectivity, give $A=\pi_0^\Delta(X)$ the discrete topology.
Exercise~\ref{ex:8} supplies compatible maps $q_n\colon X_n\to A$
which send a simplex to the common component of its vertices.
The maps
\[
 X_n\times\Delta^n_{\mathrm{top}}\longrightarrow A,
 \qquad (x,t)\longmapsto q_n(x),
\]
are continuous and respect the identifications defining realization,
because $q_m(\alpha^*x)=q_n(x)$ for every
$\alpha\colon[m]\to[n]$. They therefore descend to a continuous map
$q\colon|X|\to A$ by Definition~1.1.35.
If two realized vertices are joined by a path, their images under $q$
are joined by a path in the discrete space $A$ and thus are equal.
Their classes in $\pi_0^\Delta(X)$ coincide, proving injectivity.
The map $\phi$ is induced by vertices and commutes with simplicial
maps, so the bijection is canonical.
\end{proof}
\end{exercise}
